Use the Laplace transform to reduce the partial differential equation
\begin{equation}
\frac{\partial^2 u}{\partial x^2}
+
y\frac{\partial u}{\partial y}
=
x
\label{50000034:pde}
\end{equation}
on the domain
\[
\Omega
=
\{
(x,y)
\mid
x>0\wedge 0<y<2\pi
\}
\]
with boundary conditions
\begin{equation}
\left.
\begin{aligned}
u(x,0)                             &= -1      &\qquad &\forall x>0\\
\frac{\partial u}{\partial y}(x,0) &= e^{-x}  &\qquad &\forall x>0\\
u(0,y)                             &= \sin(y) &\qquad &0<y<2\pi   \\
\frac{\partial u}{\partial x}(0,y) &= \cos(y) &\qquad &0<y<2\pi
\end{aligned}
\quad\right\}
\label{50000034:boundary} \end{equation}
to a family of ordinary differential equations with suitable boundary
conditions.
Do the boundary conditions~\eqref{50000034:boundary} uniquely determine
the solutions of the equation~\eqref{50000034:pde}?

\begin{loesung}
The transform has to be done with respect to the $x$-coordinate,
we write the transformed coordinate as $s$.
The transform of the differential equation becomes
\begin{align*}
-
\frac{\partial u}{\partial x}(0, y)
-
su(0,y)
+
s^2\mathscr{L}u(s,y)
+
y
\frac{\partial}{\partial y} \mathscr{L}u(s, y)
&=
\frac{1}{s^2}.
\intertext{Writing $\mathscr{L}u(s,y)=U_s(y)$, this equation becomes}
-
\cos(y)
-
s\sin(y)
+
s^2 U_s(y)
+
yU'_s(y)
&=
\frac{1}{s^2}
\\
yU'_s(y)
+
s^2U_s(y)
&=
\cos(y)+s\sin(y)+\frac{1}{s^2},
\end{align*}
which is a family of linear first order ordinary differential
equations with parameter $s$.
It also satisfies the boundary conditions obtained from transforming
the first two boundary conditions in \eqref{50000034:boundary}
\begin{align*}
\mathscr{L}u(s,0)
&=
-\frac1s
&&\Rightarrow&
U_s(0)&=-\frac{1}{s}
\\
\frac{\partial \mathscr{L}u}{\partial y}(s,0)
&=
\frac{1}{s+1}
&&\Rightarrow&
U'_s(0)
&=
\frac{1}{1+s}.
\end{align*}
As the boundary conditions uniquely determine the function $U_s(y)$,
there can only be a single solution for $u(x,y)$.
\end{loesung}

\begin{diskussion}
The Laplace transform of the function $x$ may not be as well known
as the other transforms, here is its derivation:
\begin{align*}
(\mathscr{L}x)(s)
&=
\int_0^{\infty} xe^{-sx}\,dx
=
\biggl[
-\frac{xe^{-sx}}{s}
\biggr]_0^\infty
+
\int_0^\infty
\frac{e^{-sx}}{s}
\,dx
=
-\frac{1}{s}
\biggl[
\frac{e^{-sx}}s
\biggr]_0^\infty
=
\frac{1}{s^2}.
\end{align*}
\end{diskussion}

\begin{bewertung}
Choice of coordinate for the Laplace transform ({\bf C}) 1 point,
Laplace transform of derivatives ({\bf D}) 1 point,
substitution of boundary conditions ({\bf S}) 1 point,
Laplace transform of the equation ({\bf E}) 1 point,
Laplace transform of boundary conditions ({\bf B}) 1 point,
uniqueness of solution ({\bf U}) 1 point.
\end{bewertung}

